% ---
% filename : ["electrotechnique_puissance_20260429_312_mesure_puissance_compteur.tex"]
% id : "electrotechnique_puissance_20260429_312"
% title : "Mesure de puissance triphasée par compteur à disque"
% domain : "Électrotechnique"
% subdomain : "Puissance"
% tags : ["compteur d'énergie", "puissance active", "puissance apparente", "facteur de puissance", "triphasé", "énergie électrique", "ts"]
% difficulty : 2
% time_solve : 12  # Calcul : 2 (difficulte) * (3 questions * 2)
% has_octave : true
% octave_files : ["electrotechnique_puissance_20260429_312_mesure_puissance_compteur.m"]
% author : "om"
% ---
\documentclass[a4paper,11pt,answers,addpoints]{exam}
% ---------------------------------------------------------
% LANGUE ET ENCODAGE
% ---------------------------------------------------------
\usepackage[utf8]{inputenc}
\usepackage[T1]{fontenc}   % Ajout recommandé (manquait)
\usepackage[french]{babel} % Ajout pour la langue française

% ---------------------------------------------------------
% GÉOMÉTRIE ET MISE EN PAGE
% ---------------------------------------------------------
\usepackage[left=1.7cm,right=1.5cm,top=2cm,bottom=1cm]{geometry}
\usepackage{microtype}      % Meilleure répartition du texte
\usepackage{float}          % Positionnement [H]
\usepackage{needspace}      % Saut conditionnel intelligent

% Éviter les veuves/orphelines (optionnel, à décommenter si besoin)
% \widowpenalty=10000
% \clubpenalty=10000

% Espacement vertical flexible
\raggedbottom
\setlength{\parskip}{0.5em plus 0.2em minus 0.1em}
\setlength{\parindent}{0pt}

% ---------------------------------------------------------
% MATHÉMATIQUES
% ---------------------------------------------------------
\usepackage{amsmath, amssymb, bm, esvect}

% Vecteurs
\newcommand{\V}{\overrightarrow}
\def\vect#1{\overrightarrow{\strut#1}}
\providecommand{\Degrees}[0]{\ensuremath{^\circ}}

% ---------------------------------------------------------
% UNITÉS ET GRANDEURS
% ---------------------------------------------------------
\usepackage{siunitx}
\sisetup{output-decimal-marker={,}}
\DeclareSIUnit \var {var}
\DeclareSIUnit \VA {VA}
\DeclareSIUnit \kVA {kVA}
\DeclareSIUnit[quantity-product={}] \degree{\text{\textdegree}}

% ---------------------------------------------------------
% GRAPHIQUES (TikZ + CircuitikZ + PGFPlots)
% ---------------------------------------------------------
\usepackage{tikz}
\usetikzlibrary{
	arrows.meta,
	intersections,
	decorations.pathreplacing,
	positioning
}

\usepackage[european,straightvoltages,EFvoltages]{circuitikz}
\usepackage{tkz-fct}
\usepackage{pgfplots}
\pgfplotsset{compat=1.12}

% Symbole étoile-triangle personnalisé
\newcommand{\wye}{%
	\mathbin{\tikz[x=1ex,y=1ex]{%
			\draw[line width=.1ex] (0,0)--(45:1)--++(-45:1) (45:1)--++(0,1);
	}}%
}

% Sankey diagram
\usepackage{sankey}
\pgfdeclarelayer{background}
\pgfdeclarelayer{foreground}
\pgfdeclarelayer{sankeydebug}
\pgfsetlayers{background,main,foreground,sankeydebug}

% ---------------------------------------------------------
% TABLEAUX
% ---------------------------------------------------------
\usepackage{tabularx}
\usepackage{tabu}      % Alternative à tabularx (garde les deux, pas de redondance critique)
\usepackage{booktabs}  % Pour des tableaux propres

% ---------------------------------------------------------
% HYPERLIENS
% ---------------------------------------------------------
\usepackage[colorlinks=true,
menucolor=red,
breaklinks=true,
urlcolor=blue,
linkcolor=blue]{hyperref}

% ---------------------------------------------------------
% CODE (LISTINGS)
% ---------------------------------------------------------
\usepackage{xcolor}
\usepackage{listings}

\lstdefinestyle{octavestyle}{
	language=Matlab,
	basicstyle=\ttfamily\small,
	keywordstyle=\color{blue},
	commentstyle=\color{green!50!black},
	stringstyle=\color{red},
	stepnumber=1,
	frame=single,
	breaklines=true,
	breakatwhitespace=true,
	tabsize=4
}

% ---------------------------------------------------------
% MISE EN TEXTE DIVERSES
% ---------------------------------------------------------
\usepackage{enumitem}
\setlist[enumerate]{label=\Alph*)}
% ---------------------------------------------------------
% COMMANDES PERSONNELLES
% ---------------------------------------------------------
\newcommand\nomauteur{bdminasmorgul@protonmail.com}
\newcommand\dateTe{29.04.2026}
\newcommand\brancheTe{TS}
% ---------------------------------------------------------
% STYLE DE L'EXERCICE
% ---------------------------------------------------------
\pagestyle{headandfoot}
\firstpageheadrule
\runningheadrule

\firstpageheader{\brancheTe\ (\nomauteur)}{Élève : }{(\dateTe) \thepage/\numpages}
\runningheader{\brancheTe\ (\nomauteur)}{Élève : }{(\dateTe) \thepage/\numpages}

% Compteur de points en bas de page
\newcommand{\totalpointtts}{%
	\begin{tikzpicture}[remember picture, overlay]
		\node[anchor=south west, inner sep=0pt, outer sep=0pt]
		at ([xshift=0.5cm,yshift=0.5cm]current page.south west)
		{\rotatebox{90}{Total des points : \numpoints}};
	\end{tikzpicture}
}

\firstpagefooter{\totalpointtts}{}{}
\runningfooter{\totalpointtts}{}{}





% Options de l'exercice
\printanswers
\addpoints
\pointsinmargin
\bracketedpoints
\shadedsolutions
\unframedsolutions

% ---------------------------------------------------------
% DOCUMENT
% ---------------------------------------------------------
\renewcommand{\thepartno}{\arabic{partno}}
\renewcommand{\partlabel}{\thepartno)}
\begin{document}
\section*{Préparation TS PQ CFC ELMO, IELE, PELE}
\begin{questions}



\question On souhaite vérifier la puissance absorbée par une machine triphasée à l'aide d'un compteur d'énergie électromécanique. La constante du compteur est de $c = \qty{200}{[tr/kWh]}$. Lors d'une mesure, le disque du compteur effectue $n = 15$ tours en un temps $t = \qty{60}{[s]}$. 
Les mesures électriques sur la ligne indiquent une tension composée $U = \qty{400}{[V]}$ et un courant de ligne $I = \qty{8,5}{[A]}$.

\begin{parts}
	\part Déterminez l'énergie électrique $W$ consommée par la machine durant cette mesure (en \qty{}{[kWh]}).
	\part Calculez la puissance active $P$ absorbée par la machine.
	\part Calculez la puissance apparente $S$ de l'installation et déduisez-en le facteur de puissance $\cos \phi$ de la machine.
\end{parts}

\begin{solution}
	\begin{parts}
		\part L'énergie est déterminée par le nombre de rotations du disque et la constante du compteur.
		\begin{align*}
			W &= \dfrac{n}{c} = \dfrac{15}{\qty{200}{[tr/kWh]}} = \qty{0,075}{[kWh]}
		\end{align*}
		
		\part La puissance active est le rapport entre l'énergie consommée et la durée de la mesure. Il faut convertir le temps en heures : $t = \qty{60}{[s]} = \dfrac{60}{3600} = \qty{0,01667}{[h]}$.
		\begin{align*}
			P &= \dfrac{W}{t} = \dfrac{\qty{0,075}{[kWh]}}{\qty{0,01667}{[h]}} = \qty{4,5}{[kW]} = \qty{4500}{[W]}
		\end{align*}
		
		\part La puissance apparente en triphasé se calcule avec la tension composée et le courant de ligne.
		\begin{align*}
			S &= \sqrt{3} \cdot U \cdot I = \sqrt{3} \cdot \qty{400}{[V]} \cdot \qty{8,5}{[A]} = \qty{5889}{[VA]} \\[6pt]
			\cos \phi &= \dfrac{P}{S} = \dfrac{\qty{4500}{[W]}}{\qty{5889}{[VA]}} = \qty{0,764}{}
		\end{align*}
	\end{parts}
	
	\begin{verbatim}
		% Télécharger OCTAVE depuis https://wiki.octave.org/Using_Octave et l'installer
		% Puis copier-coller le code dans OCTAVE
		
		% Données du problème
		c = 200;      % Constante du compteur [tr/kWh]
		n = 15;       % Nombre de tours
		t_sec = 60;   % Temps en secondes
		U = 400;      % Tension composée [V]
		I = 8.5;      % Courant de ligne [A]
		
		% Calcul de l'énergie
		W_kwh = n / c;
		printf('Énergie mesurée W = %.4f kWh\n', W_kwh);
		
		% Calcul de la puissance active
		t_h = t_sec / 3600;
		P_watt = (W_kwh / t_h) * 1000;
		printf('Puissance active P = %.2f W\n', P_watt);
		
		% Calcul de la puissance apparente
		S_va = sqrt(3) * U * I;
		printf('Puissance apparente S = %.2f VA\n', S_va);
		
		% Facteur de puissance
		cos_phi = P_watt / S_va;
		printf('Facteur de puissance cos(phi) = %.3f\n', cos_phi);
	\end{verbatim}
\end{solution}


\end{questions}
\end{document}